\documentclass[11pt]{article}
\usepackage{fontspec}
\setmainfont{DejaVu Serif}
\usepackage{amsmath,amssymb}
\newcommand{\SUE}{\mathrm{SUE}}\newcommand{\IC}{\mathrm{IC}}\newcommand{\IR}{\mathrm{IR}}\newcommand{\AR}{\mathrm{AR}}
\usepackage{geometry}\geometry{margin=2.3cm}
\usepackage{xcolor}\usepackage{graphicx}\usepackage{enumitem}
\usepackage{booktabs}
\usepackage[hidelinks]{hyperref}

\definecolor{ink}{HTML}{16212e}\definecolor{accent}{HTML}{1f5fa8}\definecolor{soft}{HTML}{eef3fb}
\definecolor{warn}{HTML}{a8341f}
\color{ink}
\newcommand{\key}[1]{\par\medskip\noindent\colorbox{soft}{\parbox{\dimexpr\linewidth-2\fboxsep}{\textcolor{accent}{\textbf{Κλειδί:}}\ #1}}\par\medskip}
\newcommand{\ex}[1]{\par\smallskip\noindent\textbf{\textcolor{accent}{Παράδειγμα.}}\ #1\par\smallskip}
\newcommand{\warnb}[1]{\par\medskip\noindent\colorbox{soft}{\parbox{\dimexpr\linewidth-2\fboxsep}{\textcolor{warn}{\textbf{Προσοχή:}}\ #1}}\par\medskip}
\setlength{\parskip}{0.5em}\setlength{\parindent}{0pt}

\title{\textbf{Φροντιστήριο για αρχάριους: Σήματα μετοχών}\\[3pt]
\large Από την έκπληξη κερδών στο χαρτοφυλάκιο --- με απλά μαθηματικά\\[-1pt]
\normalsize\textcolor{accent}{IC $\cdot$ IR $\cdot$ t-stat $\cdot$ gate $\cdot$ συνδυασμός σημάτων $\cdot$ risk parity}}
\author{QuantIQ / AGEL AI --- Stefanos Drakos}\date{\today}

\begin{document}
\maketitle\thispagestyle{empty}

\noindent Αυτό το φυλλάδιο εξηγεί από το μηδέν, με μικρά αριθμητικά παραδείγματα, τι είναι ένα
\textbf{σήμα} (signal) στις μετοχές, πώς μετράμε αν είναι αληθινό ή τύχη, και --- το σημαντικότερο ---
\emph{πώς και γιατί ο συνδυασμός πολλών αδύναμων σημάτων} μπορεί να δώσει κάτι χρήσιμο. Δεν χρειάζεται
τίποτα πέρα από βασική άλγεβρα και την έννοια του μέσου όρου. Τα νούμερα στα παραδείγματα είναι από τα
δικά μας πραγματικά πειράματα (own-PEAD, S\&P, 2015--2024).

\tableofcontents
\bigskip

%====================================================================
\section{Τι είναι «σήμα» και γιατί είναι δύσκολο}

Στην αγορά, ό,τι φαίνεται στο διάγραμμα τιμής το βλέπουν όλοι και έχει ήδη \emph{τιμολογηθεί}. Άρα ένα
πραγματικό σήμα πρέπει να φέρνει \textbf{νέα πληροφορία} ή να εκμεταλλεύεται μια \textbf{συστηματική
ανθρώπινη συμπεριφορά} που δεν έχει διορθωθεί.

Το παράδειγμά μας είναι το \textbf{PEAD} (Post-Earnings-Announcement Drift): όταν μια εταιρεία ανακοινώνει
κέρδη \emph{καλύτερα} από το αναμενόμενο, η μετοχή της τείνει να συνεχίζει να ανεβαίνει για εβδομάδες ---
η αγορά αντιδρά \emph{αργά} (underreaction). Αυτό είναι τεκμηριωμένο από το 1968 (Ball \& Brown) και 1989
(Bernard \& Thomas).

\key{Σήμα = ένας αριθμός ανά μετοχή ανά μέρα, που \emph{προβλέπει} (έστω ελάχιστα) τη μελλοντική
σχετική απόδοση. «Σχετική» = σε σχέση με τις άλλες μετοχές, όχι απόλυτη.}

\subsection{Η «έκπληξη κερδών»: SUE}
Πώς μετράμε πόσο μεγάλη ήταν η έκπληξη; Με το \textbf{SUE} (Standardized Unexpected Earnings):
\[
\SUE_{i,t}=\frac{E_{i,t}-E_{i,t-4}}{\sigma_i},
\]
όπου $E_{i,t}$ τα κέρδη ανά μετοχή (EPS) της εταιρείας $i$ στο τρίμηνο $t$, $E_{i,t-4}$ το \emph{ίδιο
τρίμηνο πέρσι} (η «αναμονή»: seasonal random walk), και $\sigma_i$ η τυπική απόκλιση αυτών των
ετήσιων μεταβολών. Δηλαδή: «πόσες τυπικές αποκλίσεις πάνω/κάτω από το συνηθισμένο ήταν η μεταβολή».

\ex{Η εταιρεία έβγαλε EPS \$1.20 φέτος, \$1.00 πέρσι το ίδιο τρίμηνο. Η μεταβολή είναι $+\$0.20$. Αν
τα τελευταία χρόνια οι μεταβολές της είχαν τυπική απόκλιση \$0.10, τότε $\SUE=0.20/0.10=+2.0$: μια
αρκετά μεγάλη \emph{θετική} έκπληξη. Αν είχε βγάλει \$0.90 (μεταβολή $-\$0.10$), $\SUE=-1.0$.}

%====================================================================
\section{Από το σήμα στο χαρτοφυλάκιο: το long--short book}

Ένα σήμα δεν λέει «η AAPL θα ανέβει». Λέει «οι μετοχές με μεγαλύτερο SUE θα τα πάνε \emph{καλύτερα από}
αυτές με μικρότερο». Άρα το παίζουμε \textbf{σχετικά}, με πολλές μετοχές μαζί:

\begin{enumerate}[leftmargin=1.4em]
\item Κάθε μέρα, κοίτα ποιες μετοχές ανακοίνωσαν πρόσφατα και πάρε το SUE τους.
\item \textbf{Κανονικοποίηση (z-score)}: αφαίρεσε τον μέσο, διαίρεσε με την τυπική απόκλιση, ώστε όλα
τα SUE της μέρας να είναι συγκρίσιμα (μέσος 0, τυπική απόκλιση 1).
\item \textbf{Long} (αγόρασε) τις κορυφαίες θετικές, \textbf{short} (πούλα) τις χειρότερες αρνητικές.
\item \textbf{Dollar-neutral}: ίσα λεφτά long και short. Έτσι δεν στοιχηματίζεις αν η αγορά γενικά
ανεβαίνει/πέφτει --- μόνο στη \emph{σχετική} απόδοση.
\item Κράτα τη θέση όσο διαρκεί το drift (στα δεδομένα μας $\approx 51$ εργάσιμες $\approx 2.5$ μήνες),
μετά βγες. Το SUE \emph{δεν} ξαναϋπολογίζεται κάθε μέρα --- είναι ένας αριθμός ανά τριμηνιαία ανακοίνωση.
\end{enumerate}

\ex{Πραγματική μέρα (2024-07-01, 401 μετοχές): 166 long / 209 short. Κέρδισαν κυρίως τα \emph{shorts
που έπεσαν}: η BMY ήταν short και έπεσε $-1.12\%\to+4.16$ bp· η UNH short, $-1.48\%\to+4.55$ bp. Ένα
long μάς πλήγωσε: η XEL long, αλλά έπεσε $-1.05\%\to-1.44$ bp. Σύνολο μέρας: $+13.4$ bp.
($1$ bp $=0.01\%$.)}

\key{Δύο πηγές κέρδους: longs που \emph{ανεβαίνουν} \textbf{και} shorts που \emph{πέφτουν}. Κάθε μέρα
είναι μικροσκοπικό νούμερο· το αποτέλεσμα έρχεται από το άθροισμα χιλιάδων ημερών.}

%====================================================================
\section{Μετράμε το σκιλ: IC (Information Coefficient)}

Το \textbf{IC} είναι η \emph{συσχέτιση} μεταξύ του σήματος σήμερα και της απόδοσης αύριο, σε όλες τις
μετοχές μιας μέρας:
\[
\IC_d=\mathrm{corr}\big(\text{σήμα}_{i,d},\ \text{απόδοση}_{i,d+1}\big),\qquad -1\le \IC\le 1.
\]
$\IC=0$ σημαίνει «καμία προβλεπτική ικανότητα». $\IC=1$ «τέλεια πρόβλεψη» (δεν υπάρχει ποτέ).

\warnb{Στις μετοχές, ένα \emph{πολύ καλό} σήμα έχει $\IC\approx 0.05$. Τα δικά μας σήματα: $\IC\approx
0.02$--$0.04$. Φαίνεται απελπιστικά μικρό --- και όμως, με αρκετές μετοχές, αξιοποιείται. Δες παρακάτω
γιατί.}

\ex{Αν κατατάξεις 100 μετοχές με το σήμα και αύριο οι «καλές» τα πάνε λίγο καλύτερα κατά μέσο όρο, η
συσχέτιση μπορεί να είναι μόλις $0.03$. Σε μία μέρα αυτό είναι θόρυβος. Σε $2500$ μέρες $\times 300$
μετοχές, το σταθερό κομματάκι ξεχωρίζει από τον θόρυβο.}

%====================================================================
\section{Μετράμε την ποιότητα: IR (Information Ratio)}

Το \textbf{IR} είναι «απόδοση ανά μονάδα ρίσκου», ετησιοποιημένο. Αν $r_d$ η ημερήσια απόδοση του
long--short book:
\[
\IR=\sqrt{252}\,\frac{\bar r}{\sigma_r},
\]
όπου $\bar r$ ο μέσος ημερήσιος και $\sigma_r$ η τυπική απόκλιση των ημερήσιων αποδόσεων· το $\sqrt{252}$
(εργάσιμες/έτος) το κάνει ετήσιο. Είναι ο «ξάδερφος» του Sharpe ratio.

\ex{Στα δικά μας (own-PEAD, 401 μετοχές, 2015--2024): μέσος $\bar r=+0.68$ bp/μέρα, τυπική απόκλιση
$\sigma_r=11$ bp. Τότε
\[
\IR=\sqrt{252}\cdot\frac{0.68}{11}=15.87\cdot 0.062\approx 0.98\ \text{(full sample)},
\]
και $0.74$ στο αυστηρό out-of-sample. Νικηφόρες μέρες: $53.8\%$ (μέση νίκη $+7.6$ bp), χαμένες $46.1\%$
(μέση ζημιά $-7.4$ bp). Το «edge» είναι κυριολεκτικά ξυράφι: κερδίζεις \emph{λίγο} πιο συχνά και
\emph{λίγο} περισσότερο.}

\key{$\IC$ = πόσο σκιλ (μέγεθος σήματος). $\IR$ = πόσο καλά πληρώνεται αυτό το σκιλ σε σχέση με το
ρίσκο. Καλό μεμονωμένο σήμα σε free data: $\IR$ περίπου $0.3$--$0.8$.}

%====================================================================
\section{Αληθινό ή τύχη; t-statistic και Newey--West}

Ένα θετικό IR μπορεί να είναι \emph{τύχη}. Το \textbf{t-statistic} ρωτάει: «αν στην πραγματικότητα δεν
υπήρχε κανένα edge, πόσο απίθανο θα ήταν να δούμε αποτέλεσμα τόσο μεγάλο, καθαρά από τύχη;»
\[
t=\frac{\bar r}{\mathrm{SE}(\bar r)},\qquad \mathrm{SE}(\bar r)=\frac{\sigma_r}{\sqrt{n}}.
\]

\ex{\textbf{Αναλογία με κέρμα.} Ρίχνεις 10 φορές, βγαίνουν 7 κορώνες: μπορεί τύχη. Ρίχνεις 1000 φορές,
βγαίνουν 700: σχεδόν σίγουρα πειραγμένο. Το $t$ συνδυάζει \emph{πόσο μεγάλη} η απόκλιση \emph{επί πόσα
δεδομένα}: όσο περισσότερες μέρες, τόσο μικρότερη η αβεβαιότητα $\mathrm{SE}$, τόσο μεγαλύτερο το $t$.}

\begin{center}
\begin{tabular}{ccl}
\toprule
$|t|$ & Πιθανότητα να είναι τύχη & Ετυμηγορία\\
\midrule
$\sim 1$ & $\sim 30\%$ & μπορεί κάλλιστα να είναι θόρυβος\\
$\sim 2$ & $\sim 5\%$ & το κατώφλι «στατιστικά σημαντικό»\\
$2.54$ & $\sim 1\%$ & αρκετά σίγουροι ότι είναι πραγματικό\\
\bottomrule
\end{tabular}
\end{center}

\textbf{Newey--West:} τα κέρδη μαζεύονται σε «σεζόν» και οι θέσεις επικαλύπτονται, οπότε το απλό $t$
είναι ψεύτικα αισιόδοξο. Η διόρθωση Newey--West (1987) το \emph{μειώνει} ώστε να μην αυτο-εξαπατηθούμε.
Το δικό μας own-PEAD: $\IR=0.74$ με Newey--West $t=2.54$ --- δηλαδή $\sim 1\%$ πιθανότητα τύχης. Πραγματικό,
αλλά μέτριο.

%====================================================================
\section{Το «gate»: τι κρατάμε, τι πετάμε}

\key{Δεχόμαστε ένα σήμα \emph{μόνο} αν περάσει το \textbf{gate}: $|t_{\text{NW}}|>2$ out-of-sample. Όσα
βγαίνουν null ή αρνητικά \textbf{τα πετάμε} --- δεν τα «κουμπώνουμε για να γεμίσουμε».}

Από τα δικά μας:
\begin{center}
\begin{tabular}{lccl}
\toprule
Σήμα & $\IR$ & $t_{\text{NW}}$ & Απόφαση\\
\midrule
own-PEAD (ευρύ universe, mid-caps) & $0.74$ & $2.54$ & \textcolor{accent}{\textbf{κρατάμε}}\\
own-PEAD (μόνο mega-caps) & $0.40$ & $1.20$ & πετάμε (null εκεί)\\
lead-lag (ευρύ universe) & $0.03$ & $0.10$ & πετάμε (null/εύθραυστο)\\
\bottomrule
\end{tabular}
\end{center}

\warnb{Δεν μετατρέπεις σκουπίδια σε χρυσό. «Τρέχουμε πολλές αναλύσεις που βγαίνουν αρνητικές» = αυτό
είναι το \emph{κοσκίνισμα}. Οι περισσότερες υποψήφιες πέφτουν (σωστά)· κρατάς μόνο τις λίγες γνήσιες.}

%====================================================================
\section{Η καρδιά: γιατί ο συνδυασμός σημάτων βοηθάει}

Εδώ είναι το κεντρικό μάθημα. Έστω $N$ σήματα, καθένα με ίδιο μικρό $\IR$, που είναι
\textbf{ασυσχέτιστα} (λένε διαφορετικά πράγματα). Τα συνδυάζουμε ισόποσα. Τι γίνεται;

\subsection{Το μαθηματικό (απλό)}
Η απόδοση κάθε σήματος: $r_k=\mu+\varepsilon_k$, με σήμα (μέσο) $\mu$ και θόρυβο $\varepsilon_k$ τυπικής
απόκλισης $\sigma$. Ο μέσος $N$ ασυσχέτιστων:
\[
\bar r=\frac{1}{N}\sum_{k=1}^N r_k=\mu+\frac{1}{N}\sum_k \varepsilon_k.
\]
Το \textbf{σήμα μένει} $\mu$. Ο \textbf{θόρυβος} όμως, επειδή είναι ασυσχέτιστος, μαζεύεται με τυπική
απόκλιση
\[
\sigma_{\bar r}=\frac{\sigma}{\sqrt{N}}\quad(\text{όχι }\sigma).
\]
Άρα ο λόγος σήμα/θόρυβος --- δηλαδή το IR --- \emph{μεγαλώνει}:
\[
\boxed{\ \IR_{\text{combined}}=\frac{\mu}{\sigma/\sqrt{N}}=\IR_{\text{single}}\cdot\sqrt{N}\ }
\]

\ex{Δύο ασυσχέτιστα σήματα, καθένα $\IR=0.4$:
\[
\IR_{\text{combined}}=0.4\cdot\sqrt{2}=0.57.
\]
Τρία: $0.4\cdot\sqrt{3}=0.69$. Τέσσερα: $0.4\cdot\sqrt{4}=0.80$. Δέκα: $0.4\cdot\sqrt{10}=1.26$.
Ο θόρυβος πέφτει πιο γρήγορα από το σήμα --- γι' αυτό «πολλά μικρά» $>$ «ένα μεγάλο».}

\key{Είναι ο ίδιος κανόνας $\IR=\IC\cdot\sqrt{\text{breadth}}$ (Fundamental Law of Active Management,
Grinold 1989) --- αλλά εδώ το «breadth» είναι στον \emph{χώρο των σημάτων}, όχι των μετοχών.}

\subsection{Οι ΔΥΟ απαραίτητες προϋποθέσεις}
\begin{enumerate}[leftmargin=1.4em]
\item \textbf{Θετική προσδοκία} κάθε σήματος (έστω μικρή). Ένα null ($\mu\approx 0$) δεν προσθέτει
τίποτα· ένα αρνητικό βλάπτει. \emph{Γι' αυτό υπάρχει το gate.}
\item \textbf{Ορθογωνιότητα} (ασυσχέτιστα). Αν δύο σήματα είναι ίδια, ο συνδυασμός $=$ το ίδιο πράγμα.
Θες σήματα που λένε διαφορετικά πράγματα: PEAD (έκπληξη κερδών) $\perp$ momentum (τάση τιμής) $\perp$
value (αποτίμηση) $\perp$ regime (μεταβλητότητα).
\end{enumerate}

\subsection{Τι συμβαίνει αν είναι συσχετισμένα}
Αν ο συντελεστής συσχέτισης είναι $\rho$ (όχι 0), το όφελος καταρρέει. Για δύο σήματα:
\[
\IR_{\text{combined}}=\IR_{\text{single}}\cdot\sqrt{\frac{2}{1+\rho}}.
\]
\ex{Με $\rho=0$: $\sqrt{2}=1.41\times$. Με $\rho=0.5$: $\sqrt{2/1.5}=1.15\times$. Με $\rho=1$ (ίδια):
$\sqrt{2/2}=1\times$ --- κανένα όφελος. Γι' αυτό το PEAD και το lead-lag, που μοιράζονται την
earnings-season, πρέπει να διαχωριστούν (Fama--MacBeth) πριν συνδυαστούν.}

%====================================================================
\section{Πώς συνδυάζονται μηχανικά: risk parity}

Κάθε μέρα, για κάθε μετοχή, μετατρέπουμε κάθε σήμα σε z-score και τα αθροίζουμε \emph{σταθμισμένα}:
\[
\text{score}_{i}=\sum_k w_k\, z(\text{σήμα}_k)_i,\qquad w_k\propto \frac{1}{\sigma_k}\ (\text{risk parity}).
\]
Το \textbf{risk parity} δίνει \emph{λιγότερο} βάρος στο πιο θορυβώδες σήμα --- και, κρίσιμα, \emph{δεν}
χρησιμοποιεί «μαθημένα» βάρη (που θα έκαναν overfit). Μετά: long τις κορυφαίες score, short τις χειρότερες.

\ex{Σήμα A με $\sigma_A=1\%$, σήμα B με $\sigma_B=2\%$. Inverse-vol βάρη: $w_A\propto 1/1=1$,
$w_B\propto 1/2=0.5$· κανονικοποιημένα $w_A=0.67$, $w_B=0.33$. Το πιο σταθερό σήμα οδηγεί.}

\warnb{Γιατί \emph{όχι} «μάθε τα βάρη με ML/deep net»; Γιατί με $\IC\approx 0.03$ το μοντέλο θα κάνει fit
στον θόρυβο. Το DER paper το απέδειξε: η ευελιξία καταστρέφει τα σταθερά factors. Η σύσταση: low-capacity
(ridge / risk parity), όχι deep nets.}

%====================================================================
\section{Το πλήρες σύστημα και ο νόμος που το διέπει}

\[
\IR=\IC\cdot\sqrt{\text{BR}}\cdot\text{TC}
\]
($\text{BR}=$ breadth, αριθμός ανεξάρτητων στοιχημάτων· $\text{TC}=$ transfer coefficient, πόσο καθαρά
περνά το σήμα σε θέσεις μετά κόστη/περιορισμούς). Το σύστημα:
\begin{center}
\begin{tabular}{ll}
\toprule
Στρώμα & Ρόλος\\
\midrule
Signal engine & δίνει $\IC$ + breadth: PEAD, momentum, value\dots (ορθογώνια, gated)\\
Combiner & risk parity (low-capacity, χωρίς overfit)\\
Risk layer & σταθμίζει το ρίσκο ανά regime (vol targeting), προστατεύει το $\text{TC}$\\
\bottomrule
\end{tabular}
\end{center}

\ex{Με $\IC=0.03$, breadth $=300$ μετοχές $\times 4$ rebalances $=1200$, και $\text{TC}=0.3$:
$\IR=0.03\cdot\sqrt{1200}\cdot 0.3=0.03\cdot 34.6\cdot 0.3\approx 0.31$. Μέτριο --- και ρεαλιστικό.}

%====================================================================
\section{Η τίμια προειδοποίηση (μη το ξεχάσεις ποτέ)}

\begin{itemize}[leftmargin=1.4em]
\item \textbf{Κόστη.} Long--short πληρώνει spread + \emph{short financing} ($\sim 300$ bp/έτος!).
Ένα ξυράφι-edge $0.68$ bp/μέρα λεπταίνει επικίνδυνα μετά τα κόστη.
\item \textbf{Survivorship bias.} Αν το universe έχει μόνο τους «επιζώντες», φουσκώνεις το αποτέλεσμα.
Χρειάζεσαι point-in-time, as-traded μέλη.
\item \textbf{Multiple testing.} Αν δοκιμάσεις 100 σήματα, 5 θα φανούν «σημαντικά» στην τύχη. Διόρθωση:
Deflated Sharpe Ratio (Bailey \& López de Prado 2014).
\item \textbf{Fragility.} Τα σήματα αλλάζουν με το universe/περίοδο (το είδαμε: το lead-lag εξαφανίστηκε
στο ευρύ universe). Headline διαγνωστικό = \emph{durability}, όχι in-sample IR.
\item \textbf{Δεν δημιουργείται edge από το πουθενά.} Ο συνδυασμός \emph{αξιοποιεί} τα λίγα μικρά edges
που έχεις ήδη επικυρώσει --- δεν τα εφευρίσκει.
\end{itemize}

%====================================================================
\section{Σύνοψη σε μία σελίδα}
\begin{enumerate}[leftmargin=1.4em]
\item \textbf{Σήμα} = αριθμός ανά μετοχή που προβλέπει σχετική απόδοση (π.χ. SUE από έκπληξη κερδών).
\item Παίζεται \textbf{cross-sectional}, dollar-neutral long--short, με πολλές μετοχές.
\item \textbf{IC} μετράει σκιλ ($\approx 0.03$), \textbf{IR} ποιότητα ($\approx 0.3$--$0.8$),
\textbf{Newey--West t} αν είναι αληθινό ($|t|>2$).
\item \textbf{Gate}: κράτα μόνο τα $|t|>2$· πέτα τα null/αρνητικά.
\item \textbf{Συνδυασμός}: $N$ ασυσχέτιστα θετικά σήματα $\Rightarrow \IR\cdot\sqrt{N}$. Χρειάζεται
θετική προσδοκία + ορθογωνιότητα. Σταθμίζεις με risk parity (όχι ML overfit).
\item \textbf{Risk layer} από πάνω (vol targeting / DER) προστατεύει το drawdown.
\item Τίμια προσδοκία: μέτρια, εύθραυστα edges· η αξία είναι η πειθαρχημένη μεθοδολογία.
\end{enumerate}

\vfill
\noindent\rule{\linewidth}{0.4pt}\\[2pt]
\footnotesize Αναφορές: Ball \& Brown (1968)· Bernard \& Thomas (1989)· Grinold (1989, Fundamental Law)·
Newey \& West (1987)· Bailey \& López de Prado (2014, Deflated Sharpe)· Moreira \& Muir (2017,
volatility-managed). Νούμερα από QuantIQ own-PEAD, S\&P, 2015--2024 (free EDGAR data).
\end{document}
